\documentclass[runningheads]{llncs}
\usepackage[T1]{fontenc}
\usepackage{graphicx}
\usepackage{booktabs}
\usepackage{amsmath}
\usepackage{amssymb}
\usepackage{hyperref}
\usepackage{xcolor}
\usepackage{multirow}
\usepackage{subcaption}

% LNCS manual bibliography: \doi is not defined by the class; define it here.
\providecommand{\doi}[1]{\href{https://doi.org/#1}{https://doi.org/#1}}

\begin{document}

\title{Deep Learning for Cardio-Oncology Surveillance: LVEF Estimation, LV Dysfunction Classification, and Early Cardiotoxicity Prediction from Echocardiography}

\author{Shubham Innani\inst{1,2}, Michael Feldman\inst{1,2}, Spyridon Bakas\inst{1,2,3,4,\dag}}
%\author{Shubham Innani \inst{1,2} \and Suhang You\inst{1,2} \and Carla Pitarch-Abaigar\inst{1,2}, Dimitrios Makris\inst{3}, Spyridon Bakas\inst{1,2,3,4,5,\dag}}

\authorrunning{Innani et al.}

\institute{ Division of Computational Pathology, Department of Pathology and Laboratory Medicine, Indiana University School of Medicine, Indianapolis, IN, USA \and 
Indiana University Melvin and Bren Simon Comprehensive Cancer Center, USA 
\and Departments of Biostatistics and Health Data Science; Radiology and Imaging Sciences; Neurological Surgery, Indiana University School of Medicine, USA 
\and Department of Computer Science, Luddy School of Informatics, Computing, and Engineering, Indiana University, Indianapolis, IN, USA \\
    \textsuperscript{\dag} Corresponding Author: spbakas@iu.edu
}

% \authorrunning{[Author et al.]}

% \institute{[Institution TBD]\\
% \email{[email TBD]}}

\maketitle

%-----------------
\begin{abstract}
Therapy-induced cardiotoxicity is a leading non-oncological cause of morbidity in breast cancer patients receiving anthracyclines or trastuzumab, yet its subclinical precursors are routinely missed by manual echocardiographic assessment. EchoRisk-MICCAI challenge advances cardio-oncology using artificial intelligence (AI) for the early prediction and assessment of therapy-induced cardiotoxicity from echocardiography. We develop and benchmark AI models on a multicentre dataset featuring over 2,150 real-world echocardiograms across 5 longitudinal clinical timepoints address across all three EchoRisk-MICCAI 2026 tasks: (1) LVEF regression, (2) left ventricular dysfunction classification, and (3) early prediction of cardiotoxicity from pre-treatment imaging alone. A single shared architecture processes apical 4-chamber and 2-chamber cine loops through a domain-pretrained MViT-v2-S backbone (EchoPrime) with indicator-conditioned view masking, followed by concatenation-based dual-view fusion and a task-agnostic MLP head. We compare this against an R(2+1)D-18 baseline pretrained on Kinetics-400. On the public validation split, EchoPrime achieves MAE 4.24 on Task 1 (below the 5 pp clinical target), AUC-ROC 0.831 on Task 2, and AUC-ROC 0.621 on Task 3. Domain-specific pretraining outperforms the generic-video baseline on Tasks 1 and 2 and Task 3 is unresolved by backbone choice at this sample size.

\keywords{Cardiotoxicity \and Echocardiography \and LVEF Estimation \and Left Ventricular Dysfunction \and Video Understanding \and Transfer Learning}
\end{abstract}

%-----------------
\section{Introduction}

Breast cancer treatment regimens incorporating anthracyclines and trastuzumab achieve high survival rates at the cost of measurable cardiac toxicity in 20-30\% and 7-10\% of patients, respectively~\cite{ref_cardiotox_review}. Therapy-induced cardiotoxicity manifests through progressive left ventricular ejection fraction (LVEF) decline and impaired global longitudinal strain (GLS), both detectable by echocardiography before irreversible dysfunction occurs. Current clinical surveillance is manual, operator-dependent, and subject to inter-observer variability that limits sensitivity in longitudinal monitoring programmes~\cite{ref_echo_variability}. 

Automated cardiac assessment from echocardiography video advanced after EchoNet-Dynamic established LVEF estimation from apical 4-chamber cine loops at clinician-level accuracy~\cite{ref_echonet}. Subsequent work extended this to pediatric cohorts, segmentation-based pipelines, and multi-video, view-informed foundation models~\cite{ref_echonet_pediatric,ref_echoprime}. Deep learning supports clinical decision-making across medical imaging modalities, including cardiac function assessment from echocardiography video~\cite{ref_echonet,ref_echoprime}. However, no prior public challenge has paired therapy-induced cardiotoxicity labels with longitudinal echocardiographic imaging, leaving cardio-oncology surveillance without a reproducible benchmark.

The EchoRisk-MICCAI 2026 challenge addresses this gap with the first curated multicentre longitudinal echocardiography dataset with explicit cardiotoxicity labels~\cite{ref_echorisk}: 422 breast cancer patients from five European sites, 2,159 DICOM cine loops, and 1,123 clinical exams across five prospective timepoints (T1 baseline plus 3-, 6-, 9-, and 12-month follow-ups). The challenge organisers report \emph{test-set} results for a R(2+1)D-18/LSTM baseline (4.98 pp MAE on Task 1, AUC 0.849 on Task 2, AUC 0.541 on Task 3~\cite{ref_echorisk}); test labels are withheld from participants, so our results are on the public validation split and are not directly comparable. We participate in all three tasks and make the following contributions: (i) a unified dual-view fusion architecture with indicator-conditioned zero-masking for missing A2C views; (ii) an empirical comparison of a general-video backbone (R(2+1)D-18, Kinetics-400) against a domain-pretrained echocardiography backbone (EchoPrime, MViT-v2-S); and (iii) a transfer-learning study for Task 3, where the training set of 157 patients destabilises full backbone fine-tuning.

%-----------------
\section{Method}

\subsection{Problem Formulation}

All three tasks share the same dual-view DICOM input: apical 4-chamber (A4C) and apical 2-chamber (A2C) cine loops per examination. Table~\ref{tab:tasks} summarises the output format, primary evaluation metric, and supervision label for each task.

\begin{table}[t]
\caption{Task definitions for EchoRisk-MICCAI 2026.}\label{tab:tasks}
\centering
\begin{tabular}{llll}
\toprule
Task & Output & Primary metric & Label \\
\midrule
1 - LVEF Estimation & $\hat{y} \in [0,100]$ & MAE $\downarrow$ (target $\le$ 5 pp) & \texttt{lvef} \\
2 - LV Dysfunction & $\hat{p} \in [0,1]$ & AUC-ROC $\uparrow$ & \texttt{lv\_dysfunction} \\
3 - Cardiotoxicity & $\hat{p} \in [0,1]$ & AUC-ROC $\uparrow$ & \texttt{cardiotoxicity} \\
\bottomrule
\end{tabular}
\end{table}

LV dysfunction is defined per the challenge protocol as GLS $\ge -16\%$ (primary; more sensitive than LVEF for subclinical dysfunction~\cite{ref_echorisk}) or LVEF $<50\%$ (fallback when GLS is unavailable, applying to $\approx$15\% of exams), consistent with ESC 2022 cardio-oncology guidelines~\cite{ref_esc2022}. Task 3 uses only the pre-treatment (T1) examination and has a separate, smaller cohort: 280 patients in the full dataset (organiser-reported~\cite{ref_echorisk}), of which 196 have released labels (157 train / 39 validation; 84 withheld as test). The positive rate is 65.4\% across the full 280-patient cohort~\cite{ref_echorisk}, and 67.5\% (106/157) in the training split used here, reflecting enriched high-risk enrolment.

\subsection{Architecture}

Figure~\ref{fig:arch} illustrates the pipeline. A shared backbone encoder processes each view independently, producing 512-d embeddings $f_4$ (A4C) and $f_2$ (A2C). When an A2C video is absent, $f_2$ is zeroed by multiplying with the observed binary indicator \texttt{has\_a2c} $\in \{0,1\}$ (indicator-conditioned zero-masking; no learnable gating parameter). The two embeddings are concatenated to a 1024-d joint representation, which passes through a task-agnostic MLP:

\begin{equation}
  \hat{y} = W_2\,\sigma\!\left(W_1\,[f_4 \,\|\, (f_2 \cdot \texttt{has\_a2c})]\right)
\end{equation}

\noindent where $W_1 \in \mathbb{R}^{256 \times 1024}$, $W_2 \in \mathbb{R}^{1 \times 256}$, and $\sigma$ denotes ReLU followed by Dropout($p$=0.3). The raw scalar output is clamped to $[0,100]$ for Task 1 regression and passed through a sigmoid for Tasks 2-3 classification.

\begin{figure}[t]
\centering
\includegraphics[width=\textwidth]{block_diagram.png}
\caption{Dual-view echocardiography pipeline. A4C and A2C cine loops (real validation-split frames, patient ECHORISK\_0033, timepoint T1, shown) are encoded independently by a shared backbone. The A2C pathway is zeroed when that view is absent (\texttt{has\_a2c}=0). The 512-d embeddings are concatenated to 1024-d and decoded by a shared MLP head into task-specific outputs.}
\label{fig:arch}
\end{figure}

\subsection{Backbone Comparison}

We evaluate two backbones under the identical head described above.

\noindent\textbf{Baseline: R(2+1)D-18.} A factorised spatiotemporal convolution network~\cite{ref_r2plus1d} pretrained on Kinetics-400~\cite{ref_kinetics}, accessed via \texttt{torchvision}. Input shape post-channel-repeat: $(B,\,3,\,32,\,112,\,112)$ (grayscale frames repeated to 3 channels); 32 frames uniformly sampled (linspace) over the full cine loop; ImageNet normalisation.

\noindent\textbf{Primary: EchoPrime.} We load the full EchoPrime video encoder~\cite{ref_echoprime}, a multiscale vision transformer (MViT-v2-S architecture~\cite{ref_mvit}) whose final projection to 512-d is itself pretrained via contrastive vision-language alignment on over 12 million paired echocardiography video-report studies. Concretely, we instantiate the MViT-v2-S architecture with its head reshaped to \texttt{Linear(768, 512)} to match the released checkpoint, then load EchoPrime's weights in full (including the 512-d projection); no head component is randomly initialised, and no ImageNet weights are used. We use the visual encoder only; no text/retrieval components are used. Input shape: $(B,\,3,\,16,\,224,\,224)$; 16 frames uniformly sampled (\texttt{np.linspace}) over the full cine loop; normalisation statistics derived from echocardiography data (mean\,=\,[29.11,\,28.08,\,29.10], std\,=\,[47.99,\,46.46,\,47.20] on the raw 0-255 pixel scale). This deviates from EchoPrime's original pretraining configuration (32 frames at fixed stride\,2); we disclose this as a methodological difference that may contribute to the Task\,2 gap versus the organiser baseline (\S\ref{sec:discussion}).

\subsection{Task-Specific Training Strategy}

\noindent\textbf{Task 1} (both backbones) minimises Huber loss ($\delta$=5.0, robust to rare low-LVEF outliers) with full end-to-end fine-tuning of backbone and head.

\noindent\textbf{Task 2} also fine-tunes end-to-end, but the two backbones use different loss functions: the baseline uses focal loss ($\alpha$=0.75, $\gamma$=2.0) to address class imbalance, while EchoPrime uses weighted binary cross-entropy with $\texttt{pos\_weight} = n_\text{neg}/n_\text{pos}$ computed from training label counts. Both choices were tuned independently per backbone rather than held fixed, so this is a confound we note rather than control for.

\noindent\textbf{Task 3} (157 training patients, 39 validation) requires regularised transfer learning, because at this sample size unfreezing the backbone destabilises training rather than improving it, for both backbones. The baseline initialises from the Task 1 checkpoint with a fully frozen backbone (head-only training). For EchoPrime, we initialise from the Task 2 checkpoint (richer, classification-aligned features) and compare three regimes: (i) fully frozen backbone, head-only training; (ii) last 4 MViT blocks unfrozen at $\text{lr}=10^{-5}$; and (iii) last 2 MViT blocks unfrozen at $\text{lr}=10^{-6}$ with the head at $\text{lr}=10^{-4}$ (differential learning rates, single AdamW optimiser). Regime (ii) degrades validation AUC monotonically from epoch 1, indicating that unfreezing $\sim$17M parameters on 157 samples destroys pretrained features faster than the head can exploit the newly available capacity. Regime (iii) yields the best EchoPrime result and is reported below. Both backbones use weighted BCE with $\texttt{pos\_weight}$=0.481 (train set: 106 positive, 51 negative) and weight decay $10^{-3}$ (stronger than Tasks 1-2's $10^{-4}$, to counter overfitting at this sample size).

Data augmentation for Task 3 is stronger than for Tasks 1-2 to further counteract overfitting: random rotation ($\pm 12^\circ$), zoom-crop (80-100\%), Gaussian noise ($\sigma \in [3,8]$, simulating speckle), and temporal reversal ($p$=0.3; used as a label-preserving augmentation, not as a claim of physiological equivalence between systolic and diastolic ordering).

%-----------------
\section{Experimental Setup}

\subsection{Dataset and Splits}

The EchoRisk-MICCAI 2026 full release~\cite{ref_echorisk} contains 422 breast cancer patients from five European sites, 2,159 DICOM cine loops, and 1,123 clinical exams across five longitudinal timepoints (T1 baseline plus four 3-month follow-ups), split 56\% train / 14\% validation / 30\% test (test labels withheld by organisers). The dataset covers all image quality levels and includes multi-vendor hardware (GE, Philips, Siemens, Canon), reflecting real-world acquisition variability; 92\% of exams in Tasks 1-2 have dual-view coverage~\cite{ref_echorisk}. The publicly released train and validation portion used in this work, counted directly from the released label files, comprises: for Tasks 1-2, 296 patients and 786 exams across all 5 timepoints (630 train / 156 validation; 775 A4C and 728 A2C videos, 1,503 videos total); for Task 3 (baseline T1 only), 196 patients and 196 exams (157 train / 39 validation; 196 A4C and 196 A2C videos, 392 videos total).

\subsection{Preprocessing}

DICOMs are \texttt{YBR\_FULL\_422} JPEG-compressed cine loops with high frame-count variability (median 60, IQR 44-156, range 14-711~\cite{ref_echorisk}). The Y (luma) channel is extracted as grayscale, uniformly sampled via \texttt{np.linspace} to 32 frames (baseline) or 16 frames (EchoPrime), and resized to $112{\times}112$ or $224{\times}224$, respectively. No cardiac segmentation or motion-based frame selection is applied; the backbone learns temporal saliency from raw pixel sequences.

\subsection{Hyperparameters}

All runs use AdamW with weight decay $10^{-4}$ (Tasks 1-2) or $10^{-3}$ (Task 3, both backbones). Batch size is 4. The cosine annealing schedule is applied across 50 epochs for Tasks 1-2 and 30 (baseline) / 40 (EchoPrime) epochs for Task 3. EchoPrime runs use automatic mixed precision (AMP) with gradient norm clipping at 1.0. Model selection uses the checkpoint with the best validation metric (lowest MAE for Task 1, highest AUC for Tasks 2-3) evaluated once per epoch. Reported validation metrics therefore reflect best-epoch selection on the same split, introducing optimistic bias especially for Task 3 ($n$=39); we defer cross-validation to the camera-ready submission.

%-----------------

%-----------------
\section{Results}

All metrics are computed with the official \texttt{echorisk-benchmark} scoring package on the public validation split.

\begin{table}[t]
\caption{Task 1 - LVEF estimation (MAE $\downarrow$). $n=156$ validation exams. Clinical target: MAE $\le$ 5 pp. Bold: best per column among our two models.}\label{tab:task1}
\centering
\begin{tabular}{llll}
\toprule
Backbone & MAE & RMSE & Pearson $r$ \\
\midrule
R(2+1)D-18 (Kinetics-400) & 4.622 & 5.968 & 0.524 \\
\textbf{EchoPrime (MViT-v2-S)} & \textbf{4.243} & \textbf{5.531} & \textbf{0.625} \\
\bottomrule
\end{tabular}
\end{table}

\begin{table}[t]
\caption{Task 2 - LV dysfunction classification (AUC-ROC $\uparrow$). $n=156$ validation exams. Bold: best per column among our two models. Both models reach the same balanced accuracy at the fixed 0.5 threshold via different sensitivity/specificity trade-offs (a coincidence at this sample size), despite differing AUC. Organiser R(2+1)D-18/LSTM baseline: AUC 0.849 on held-out test set~\cite{ref_echorisk} (not directly comparable).}\label{tab:task2}
\centering
\begin{tabular}{llll}
\toprule
Backbone & AUC-ROC & Balanced Acc. & Sens.@90\%\,Spec. \\
\midrule
R(2+1)D-18 (Kinetics-400) & 0.762 & 0.678 & 0.292 \\
\textbf{EchoPrime (MViT-v2-S)} & \textbf{0.831} & \textbf{0.678} & \textbf{0.458} \\
\bottomrule
\end{tabular}
\end{table}

\begin{table}[t]
\caption{Task 3 - Early cardiotoxicity prediction (AUC-ROC $\uparrow$). $n=39$ validation patients. Clinical target: sensitivity $\ge$80\% at FPR 10-20\%. Bold within each metric column marks the best result; note that the top-AUC model (R(2+1)D-18, 0.636) achieves balanced accuracy 0.500, i.e.\ chance-level classification despite the higher AUC.}\label{tab:task3}
\centering
\begin{tabular}{lllll}
\toprule
Backbone / regime & AUC-ROC & Bal.\,Acc. & Sens.@FPR10 & Brier \\
\midrule
R(2+1)D-18, frozen (Task-1-init) & \textbf{0.636} & 0.500 & 0.250 & 0.250 \\
EchoPrime, frozen head-only (T2-init) & 0.618 & - & - & - \\
EchoPrime, last-4-blocks unfrozen & 0.497 & - & - & - \\
EchoPrime, last-2-blocks, lr=$10^{-6}$ & 0.621 & \textbf{0.555} & \textbf{0.250} & \textbf{0.255} \\
\bottomrule
\end{tabular}
\end{table}

Figure~\ref{fig:results} visualises these results directly on the validation set: per-exam predicted-vs-true LVEF (Task 1) and ROC curves (Tasks 2-3).

\begin{figure}[t]
\centering
\begin{subfigure}[b]{0.32\textwidth}
\includegraphics[width=\textwidth]{fig2a_task1_scatter.png}
\caption{Task 1: LVEF}\label{fig:res_a}
\end{subfigure}
\hfill
\begin{subfigure}[b]{0.32\textwidth}
\includegraphics[width=\textwidth]{fig2b_task2_roc.png}
\caption{Task 2: ROC}\label{fig:res_b}
\end{subfigure}
\hfill
\begin{subfigure}[b]{0.32\textwidth}
\includegraphics[width=\textwidth]{fig2c_task3_roc.png}
\caption{Task 3: ROC}\label{fig:res_c}
\end{subfigure}
\caption{Validation-set results. (a) Predicted vs. true LVEF for both backbones (dashed line: identity). (b--c) ROC curves for Task 2 and Task 3; the dotted vertical line in (c) marks the FPR=10\% clinical operating point referenced by the Task 3 clinical target.}
\label{fig:results}
\end{figure}

%-----------------
\section{Discussion}
\label{sec:discussion}

In this article, we introduced a single dual-view fusion architecture that addresses all three EchoRisk-MICCAI 2026 tasks, pairing a domain-pretrained echocardiography backbone (EchoPrime, MViT-v2-S) with indicator-conditioned view masking and a task-agnostic MLP head, and we compared it against a Kinetics-pretrained R(2+1)D-18 baseline under the same fusion and head. We validated both systems on the public validation split using the official \texttt{echorisk-benchmark} scorer, and additionally studied frozen versus partially unfrozen transfer regimes for the low-sample-size Task 3. Our results show that the echocardiography-pretrained system meets the clinical MAE target on LVEF estimation (4.24 pp) and achieves strong discrimination on LV dysfunction classification (AUC-ROC 0.831), while early cardiotoxicity prediction from pre-treatment imaging alone (Task 3) remains unresolved regardless of backbone at this sample size. We discuss each of these findings, and their limitations, below.

\noindent\textbf{Domain pretraining drives gains over the generic-video baseline on Tasks 1 and 2.} Replacing the Kinetics-pretrained R(2+1)D-18 with EchoPrime, holding the head and fusion strategy constant, reduces Task 1 MAE by 0.38 points (4.62 to 4.24) and raises Task 2 AUC-ROC by 0.069 (0.762 to 0.831). These gains are consistent with cardiac structure and motion cues captured during echocardiography-specific contrastive pretraining~\cite{ref_echoprime}. The two pipelines differ not only in pretraining but also in spatial resolution, frame count, loss function, and optimisation; this comparison should therefore be read as a comparison of two complete systems rather than as isolated evidence that domain-specific pretraining alone accounts for the observed differences. Both single-clip models fall below the organiser's R(2+1)D-18/LSTM baseline (AUC 0.849 on test); that baseline adds an LSTM temporal aggregator absent from our MLP head, which likely accounts for part of the Task 2 gap. The 16-frame / linspace sampling deviating from EchoPrime's native 32-frame / stride-2 configuration (\S2.3) is a second candidate, as it reduces the temporal resolution available to the backbone. Task 1 MAE 4.243 meets the $\le$5 pp clinical target, but Pearson $r$=0.625 indicates predictions with limited spread: a predictor regressing toward the cohort mean can meet a MAE target without discriminative ability. We report the validation EF standard deviation and a mean-predictor MAE baseline in the camera-ready submission.

\noindent\textbf{Task 3 is the primary unresolved challenge.} With 157 training patients, the EchoPrime backbone's advantage on current-state assessment does not transfer cleanly to the prognostic task. Head-only training caps at AUC 0.618, marginally below the R(2+1)D-18 baseline (0.636); that baseline itself achieves balanced accuracy 0.500 (chance-level classification), so the higher AUC reflects a ranking benefit that does not translate to useful clinical thresholding. These findings are directionally consistent with the organiser report~\cite{ref_echorisk}, which shows that the strongest video model achieves AUC 0.541 on the held-out test set (not available to us), statistically indistinguishable from a logistic-regression clinical floor, and confirms that the same architecture yields AUC 0.49-0.54 regardless of pretraining in pilot experiments. The bottleneck is predictive signal availability rather than model capacity. Handcrafted radiomics features from the same cohort show statistically significant association with cardiotoxicity outcomes~\cite{ref_echorisk}, indicating the signal exists in the imaging data but is not recovered by end-to-end video architectures at this training scale.

\noindent\textbf{Clinical gap in Task 3.} The target of sensitivity $\ge$80\% at FPR 10-20\% is unmet by any regime (Sens.@FPR10 = 0.25; Fig.~\ref{fig:res_c}). Directions unexplored here include: (i) joint multi-task training with a shared Task-1/2 backbone and a simultaneous Task-3 head; (ii) clinical feature fusion incorporating baseline LVEF and GLS as auxiliary inputs; and (iii) patient-level longitudinal modelling across available timepoints.

\noindent\textbf{Checkpoint provenance and patient-level separation (Task 3).} The challenge releases Task 1--2 and Task 3 as independently split cohorts, and their patient assignments are not aligned: the Task 1/2 training split and the Task 3 validation split are not disjoint at the patient level. Concretely, 20 of the 39 Task 3 validation patients also appear in the Task 1/2 training data, and for 19 of them the exact baseline (T1) examination used at Task 3 validation was present when training the Task 2 checkpoint from which our Task 3 model is initialised; several of these patients additionally contribute follow-up timepoints (T2--T5) to Task 1/2 training. The reported EchoPrime Task 3 model carries these representations forward (initialised from the Task 2 checkpoint, with the last two MViT blocks subsequently adapted), so its backbone has seen a subset of the Task 3 validation patients during Task 1/2 training. This means the reported Task 3 validation AUC is, if anything, optimistic relative to true held-out generalisation, which reinforces rather than weakens our central finding that Task 3 remains unresolved at this sample size. A fully patient-disjoint transfer protocol would require re-splitting the challenge cohorts and is deferred to future work; the held-out organiser test set, which we cannot self-evaluate, is not subject to this overlap.

\noindent\textbf{Other limitations.} Validation sets are small ($n=156$ for Tasks 1-2, $n=39$ for Task 3), introducing non-trivial metric variance; reported figures reflect best-epoch selection on the same validation split and carry optimistic bias, particularly for Task 3. Missing A2C views are handled via zero-masking rather than imputation. The Task 2 loss-function difference (focal vs.\ weighted BCE, chosen independently per backbone) means the Task 2 comparison is not a pure backbone ablation. Confidence intervals and cross-validation are deferred to the camera-ready submission.

%-----------------
\section{Conclusion}

A shared dual-view fusion architecture with an echocardiography-pretrained MViT-v2-S backbone (EchoPrime) outperforms a Kinetics-pretrained R(2+1)D-18 baseline on Tasks 1 and 2 on the validation split, meeting the Task 1 clinical MAE target. Neither single-clip model reaches the LSTM-augmented organiser reference on Task 2, pointing to temporal sequence modelling as an important complementary direction. Early cardiotoxicity prediction from pre-treatment imaging alone (Task 3) remains an open problem that is not resolved by backbone choice at this sample size.

%-----------------
\begin{credits}
\subsubsection{\ackname}
We thank the EchoRisk-MICCAI organising committee for curating and releasing this benchmark.

\subsubsection{\discintname}
The authors have no competing interests to declare that are relevant to the content of this article.
\end{credits}

%-----------------
\begin{thebibliography}{10}

\bibitem{ref_cardiotox_review}
Zamorano, J.L., Lancellotti, P., Rodriguez Mu\~noz, D., et al.: 2016 ESC Position Paper on cancer treatments and cardiovascular toxicity. Eur. Heart J. \textbf{37}(36), 2768-2801 (2016). \doi{10.1093/eurheartj/ehw211}

\bibitem{ref_echo_variability}
Thavendiranathan, P., Grant, A.D., Negishi, T., et al.: Reproducibility of echocardiographic techniques for sequential assessment of left ventricular ejection fraction and volumes. J. Am. Coll. Cardiol. \textbf{61}(1), 77-84 (2013). \doi{10.1016/j.jacc.2012.09.035}

\bibitem{ref_echonet}
Ouyang, D., He, B., Ghorbani, A., et al.: Video-based AI for beat-to-beat assessment of cardiac function. Nature \textbf{580}, 252-256 (2020). \doi{10.1038/s41586-020-2145-8}

\bibitem{ref_echonet_pediatric}
Reddy, C.D., Lopez, L., Ouyang, D., Zou, J.Y., He, B.: Video-based deep learning for automated assessment of left ventricular ejection fraction in pediatric patients. J. Am. Soc. Echocardiogr. \textbf{36}(5), 482-489 (2023). \doi{10.1016/j.echo.2023.01.015}

\bibitem{ref_echoprime}
Vukadinovic, M., Tang, X., Yuan, N., Cheng, P., Li, D., Cheng, S., He, B., Ouyang, D.: EchoPrime: A Multi-Video View-Informed Vision-Language Model for Comprehensive Echocardiography Interpretation. arXiv:2410.09704 (2024)

\bibitem{ref_echorisk}
Kalliatakis, G., Karanasiou, G., Manikis, G., Tsiknakis, M., Fotiadis, D., et al.: EchoRisk: A Multicentre Echocardiography Dataset and Benchmark for Cardio-Oncology. arXiv:2607.01039 (2026). \url{https://echorisk-miccai.github.io/}

\bibitem{ref_esc2022}
Lyon, A.R., L{\'o}pez-Fern{\'a}ndez, T., Couch, L.S., et al.: 2022 ESC Guidelines on Cardio-Oncology Developed in Collaboration with the European Hematology Association (EHA), the European Society for Therapeutic Radiology and Oncology (ESTRO) and the International Cardio-Oncology Society (IC-OS). Eur. Heart J. Cardiovasc. Imaging \textbf{23}(10), e333-e465 (2022). \doi{10.1093/ehjci/jeac106}

\bibitem{ref_r2plus1d}
Tran, D., Wang, H., Torresani, L., Ray, J., LeCun, Y., Paluri, M.: A Closer Look at Spatiotemporal Convolutions for Action Recognition. In: CVPR, pp. 6450-6459 (2018). \doi{10.1109/CVPR.2018.00675}

\bibitem{ref_kinetics}
Kay, W., Carreira, J., Simonyan, K., et al.: The Kinetics Human Action Video Dataset. arXiv:1705.06950 (2017)

\bibitem{ref_mvit}
Li, Y., Wu, C.Y., Fan, H., Mangalam, K., Xiong, B., Malik, J., Feichtenhofer, C.: MViTv2: Improved Multiscale Vision Transformers for Classification and Detection. In: CVPR, pp. 4804-4814 (2022). \doi{10.1109/CVPR52688.2022.00476}

\end{thebibliography}

\end{document}
